% Adapted from the TeXample global-nodes example; CC BY-SA 4.0.
\documentclass[tikz,border=8pt]{standalone}
\usepackage{minimalist-profile}
\usetikzlibrary{calc}
\begin{document}
\begin{tikzpicture}[ms base]
  \node[anchor=base west,inner sep=0bp] (base) at (0,0)
    {$\vec a_p=\vec a_o+\frac{{}^bd^2}{dt^2}\vec r+$};
  \node[anchor=base west,text=msColor5]
    (coriolis) at ([xshift=\msSpaceLabel]base.base east)
    {$2\vec\omega_{ib}\times\frac{{}^bd}{dt}\vec r$};
  \node[anchor=base west,inner sep=0bp] (plus1)
    at ([xshift=\msSpaceLabel]coriolis.base east) {$+$};
  \node[anchor=base west,text=msColor1]
    (transversal) at ([xshift=\msSpaceLabel]plus1.base east)
    {$\vec\alpha_{ib}\times\vec r$};
  \node[anchor=base west,inner sep=0bp] (plus2)
    at ([xshift=\msSpaceLabel]transversal.base east) {$+$};
  \node[anchor=base west,text=msColor6]
    (centripetal) at ([xshift=\msSpaceLabel]plus2.base east)
    {$\vec\omega_{ib}\times(\vec\omega_{ib}\times\vec r)$};
  \node[ms body,text=msColor5,above=\msSpacePanel of coriolis] (c-label)
    {Coriolis acceleration};
  \node[ms body,text=msColor1,below=\msSpacePanel of transversal] (t-label)
    {Transversal acceleration};
  \node[ms body,text=msColor6,above=\msSpacePanel of centripetal] (p-label)
    {Centripetal acceleration};
  \draw[ms arrow,draw=msColor5] (c-label.south) -- (coriolis.north);
  \draw[ms arrow,draw=msColor1] (t-label.north) -- (transversal.south);
  \draw[ms arrow,draw=msColor6] (p-label.south) -- (centripetal.north);
\end{tikzpicture}
\end{document}
